\section{Closed-loop task quality}
\label{sec:transfer-gap}

\textbf{Exp\#8: Matched-retention answer quality.}
The policy evaluated here is the runtime reconstruction of
\S\ref{sec:realization}, not the offline predictor of Table~\ref{tab:v2}. The main
comparison uses chat-formatted LongBench QA~\cite{longbench}: its seven English
question-answering datasets (NarrativeQA, Qasper, MultiFieldQA-en, HotpotQA,
2WikiMultihopQA, MuSiQue, TriviaQA; abbreviated in Fig.~\ref{fig:forest}), two
architectures (Llama-3.1-8B and Qwen2.5-7B) and 448 source prompts, 64 per dataset, each
evaluated on both architectures (896 model--prompt evaluations, 64 per
architecture--dataset cell), at 20\% logical retention, each policy driving its own
masked decoding trajectory.

\textbf{Inferential unit and margin.} Because the two architectures answer the same
questions, the 14 cells are not independent. The primary analysis averages the two
models within each dataset and runs paired TOST over the seven datasets (df 6); the
14 cells treated as independent are a sensitivity analysis. The margin is two
absolute F1 points, 5.4\% of the rerun baseline 0.372: a practical tolerance carried
over from the archived analysis, not pre-registered and not a deployment threshold.
Fig.~\ref{fig:forest} exposes task heterogeneity: mean equivalence is not per-task.

\begin{figure}[t]
\centering
\includegraphics[width=\linewidth]{fig_cell_forest.pdf}
\caption{Exp\#8, pinned bfloat16 rerun. Reconstructed-minus-H2O-style F1 per
architecture $\times$ dataset cell (64 paired requests, 95\% t intervals; the two models
of a dataset adjacent) and pooled estimates (90\% t intervals): seven datasets (primary,
filled) and 14 cells (sensitivity, hollow); the band is the $\pm0.02$ margin.}
\label{fig:forest}
\end{figure}

\begin{table}[t]
\centering
\caption{Exp\#8: masked-loop LongBench F1 at 20\% retention, 896 model--prompt
evaluations. $\Delta$ is the difference from the H2O-style accumulator with a 90\% $t$
interval over the seven datasets (df 6); $p_{\mathrm{TOST}}$ is the TOST $p$ at
$\pm.02$ over the same datasets. Pinned bfloat16 rerun, official scorer, all
references.}
\label{tab:perlayer}
\footnotesize
\setlength{\tabcolsep}{3pt}
\begin{tabular*}{\columnwidth}{@{\extracolsep{\fill}}lccc@{}}
\toprule
policy realization & F1 & $\Delta$ [90\% CI] & $p_{\mathrm{TOST}}$ \\
\midrule
%%PERLAYER-ROWS-BEGIN%%
full cache & .419 & +.047 & --- \\
H2O-style & .372 & --- & --- \\
reconstructed & .374 & +.002 [-.006, +.010] & .0023 \\
Ada-KV-style & .363 & -.008 [-.016, -.000] & .015 \\
PyramidKV-style & .349 & -.023 [-.046, +.001] & .582 \\
%%PERLAYER-ROWS-END%%
\bottomrule
\end{tabular*}
\end{table}

\textbf{Pinned rerun.} With the official LongBench scorer and all references, the
reconstructed scorer's mean F1 difference from the accumulator over 896 evaluations is
$+0.0017$: the seven-dataset TOST gives $p=0.0023$ at $\pm0.02$ and the 14-cell
sensitivity analysis $p=0.0007$. The seven-dataset $t$ interval is $[-0.0064,0.0098]$;
the 14-cell $t$-based and cell-bootstrap 90\% intervals are $[-0.0063,0.0097]$ and
$[-0.0056,0.0087]$. Supplementary $\pm0.01$ tests ($p=0.047$ seven datasets, 0.045 14 cells) use a margin
that was not set in advance, so we do not count them. Per architecture the
difference is $+0.0052$ (Llama, $p=0.016$) and $-0.0018$ (Qwen, $p=0.026$). The
Ada-KV-style comparator remains equivalent ($-0.0084$, $p=0.015$); the PyramidKV-style
one does not ($-0.023$, $p=0.58$). Both allocation comparators only move a fixed total
budget across layers (Ada-KV-style: by attention concentration; PyramidKV-style: a
linear pyramid) and select within a layer as H2O-style does, so no finding covers original H2O, Ada-KV or
PyramidKV, or equivalence between the reconstructed and Ada-KV-style policies; and this
is equivalence between compressed policies, not absence of compression loss: both stay
0.045--0.047 below full cache. The policies also differ in EMA window and mandatory
blocks (Table~\ref{tab:protocol}): this compares two complete runtime realizations, not
a scorer-only intervention, which is untested. Logged retained blocks per decode step
average 25.02 for the H2O-style, reconstructed and Ada-KV-style policies and 25.01 for
the PyramidKV-style one (median 26, range 7--26): retention is matched in counts, not
only in the nominal 0.2.

\textbf{What the rerun repaired.} Archived Qwen float16 decoding emitted all-NaN logit
rows, producing token~0 (``!''): a 56-prompt replay found this in 54 requests, against
none in bfloat16/float32. In the rerun no Qwen
full-cache output contains a run of ``!'' (0/448 versus 181/448) and 34/448 reach the
48-token cap (versus 210/448). The archived F1 difference was $+0.00511$ under a
single-reference scorer ($+0.0030$ after all-reference rescoring); the artifact keeps
the original scores and the repair record.

\textbf{Exp\#9: Archived budget and context checks.} Sweeping retention over
$\{0.10,0.20,0.30,0.50\}$ yields no resolved per-budget quality gain; at 16K the larger
Llama sample (448 prompts, 20\% retention) has mean F1 difference $-0.0006$,
request-level 95\% interval $[-0.012,0.012]$, and a 160-prompt pilot that had suggested
$+0.010$ is not sustained, so we claim no long-context gain. A synthetic needle
probe is a negative control, not evidence about query-aware selection. Neither check
proves curve equivalence or zero memory saving.
